% Source: https://artofproblemsolving.com/community/c3h2535947_sum_confirmation_plz
% Stable identity: aops-topic-2535947-q01
\textbf{(a)} All terms from $3$ to $N$ cancel, leaving $1+2-1=\boxed{2}$ for every $N\ge3$.

\textbf{(b)} Now the terms $3,\dots,N$ cancel and the terms $N+1,\dots,2N$ remain:
\[1+2-1+\sum_{n=N+1}^{2N}n=2+\dfrac{N(3N+1)}{2},\]
which tends to infinity as $N\to\infty$.

\textbf{(c)} An infinite sum $\sum a_n$ means the limit of its partial sums. The partial sums $1+2+\cdots+N=\dfrac{N(N+1)}2$ grow without bound, so $\sum_{n=1}^{\infty}n$ does not exist as a number (it diverges), and neither does $\sum_{n=3}^{\infty}n$. Subtracting two divergent series is not defined: as (a) and (b) show, cutting off the two sums at different places gives completely different limits ($2$, or $\infty$; other cut-offs give any other behaviour). So the expression has no value. The finite statement in (a) is correct.
